\section{Modelo Fisico}

\subsection{Implementacao Completa}

\subsubsection{Tabelas Complementares}

\paragraph{Tabela: orcamentos}
\begin{lstlisting}[caption=Estrutura da tabela orcamentos]
CREATE TABLE orcamentos (
    id UUID PRIMARY KEY DEFAULT gen_random_uuid(),
    nome VARCHAR(100) NOT NULL,
    descricao TEXT,
    valor_planejado DECIMAL(15,2) NOT NULL,
    valor_gasto DECIMAL(15,2) DEFAULT 0.00,
    periodo_inicio DATE NOT NULL,
    periodo_fim DATE NOT NULL,
    categoria_id UUID REFERENCES categorias(id) ON DELETE CASCADE,
    usuario_id UUID REFERENCES usuarios(id) ON DELETE CASCADE,
    grupo_id UUID REFERENCES grupos(id) ON DELETE CASCADE,
    tipo_periodo VARCHAR(20) DEFAULT 'mensal',
    alerta_percentual INTEGER DEFAULT 80,
    ativo BOOLEAN DEFAULT true,
    created_at TIMESTAMP DEFAULT CURRENT_TIMESTAMP,
    updated_at TIMESTAMP DEFAULT CURRENT_TIMESTAMP,
    
    CONSTRAINT chk_orcamentos_periodo CHECK (periodo_inicio <= periodo_fim),
    CONSTRAINT chk_orcamentos_valor CHECK (valor_planejado > 0),
    CONSTRAINT chk_orcamentos_alerta CHECK (alerta_percentual BETWEEN 0 AND 100),
    CONSTRAINT chk_orcamentos_escopo CHECK (
        (usuario_id IS NOT NULL AND grupo_id IS NULL) OR
        (usuario_id IS NULL AND grupo_id IS NOT NULL)
    )
);
\end{lstlisting}

\paragraph{Tabela: metas}
\begin{lstlisting}[caption=Estrutura da tabela metas]
CREATE TABLE metas (
    id UUID PRIMARY KEY DEFAULT gen_random_uuid(),
    nome VARCHAR(100) NOT NULL,
    descricao TEXT,
    valor_alvo DECIMAL(15,2) NOT NULL,
    valor_atual DECIMAL(15,2) DEFAULT 0.00,
    data_alvo DATE,
    tipo VARCHAR(30) NOT NULL,
    categoria_id UUID REFERENCES categorias(id) ON DELETE SET NULL,
    conta_destino_id UUID REFERENCES contas(id) ON DELETE SET NULL,
    usuario_id UUID REFERENCES usuarios(id) ON DELETE CASCADE,
    grupo_id UUID REFERENCES grupos(id) ON DELETE CASCADE,
    frequencia_contribuicao VARCHAR(20),
    valor_contribuicao DECIMAL(15,2),
    ativa BOOLEAN DEFAULT true,
    alcancada BOOLEAN DEFAULT false,
    data_alcancada TIMESTAMP,
    created_at TIMESTAMP DEFAULT CURRENT_TIMESTAMP,
    updated_at TIMESTAMP DEFAULT CURRENT_TIMESTAMP,
    
    CONSTRAINT chk_metas_valor CHECK (valor_alvo > 0),
    CONSTRAINT chk_metas_tipo CHECK (
        tipo IN ('poupanca', 'investimento', 'emergencia', 'compra', 'viagem', 'outros')
    ),
    CONSTRAINT chk_metas_frequencia CHECK (
        frequencia_contribuicao IS NULL OR 
        frequencia_contribuicao IN ('diario', 'semanal', 'mensal', 'anual')
    ),
    CONSTRAINT chk_metas_escopo CHECK (
        (usuario_id IS NOT NULL AND grupo_id IS NULL) OR
        (usuario_id IS NULL AND grupo_id IS NOT NULL)
    )
);
\end{lstlisting}

\paragraph{Tabela: recorrencias}
\begin{lstlisting}[caption=Estrutura da tabela recorrencias]
CREATE TABLE recorrencias (
    id UUID PRIMARY KEY DEFAULT gen_random_uuid(),
    nome VARCHAR(100) NOT NULL,
    descricao TEXT,
    valor DECIMAL(15,2) NOT NULL,
    tipo VARCHAR(20) NOT NULL,
    frequencia VARCHAR(20) NOT NULL,
    dia_vencimento INTEGER,
    data_inicio DATE NOT NULL,
    data_fim DATE,
    conta_id UUID NOT NULL REFERENCES contas(id) ON DELETE CASCADE,
    categoria_id UUID NOT NULL REFERENCES categorias(id) ON DELETE RESTRICT,
    usuario_id UUID NOT NULL REFERENCES usuarios(id) ON DELETE CASCADE,
    grupo_id UUID REFERENCES grupos(id) ON DELETE SET NULL,
    proxima_execucao DATE NOT NULL,
    ativa BOOLEAN DEFAULT true,
    auto_confirmar BOOLEAN DEFAULT false,
    observacoes TEXT,
    created_at TIMESTAMP DEFAULT CURRENT_TIMESTAMP,
    updated_at TIMESTAMP DEFAULT CURRENT_TIMESTAMP,
    
    CONSTRAINT chk_recorrencias_tipo CHECK (
        tipo IN ('receita', 'despesa')
    ),
    CONSTRAINT chk_recorrencias_frequencia CHECK (
        frequencia IN ('diario', 'semanal', 'quinzenal', 'mensal', 'bimestral', 
                      'trimestral', 'semestral', 'anual')
    ),
    CONSTRAINT chk_recorrencias_dia CHECK (
        dia_vencimento IS NULL OR dia_vencimento BETWEEN 1 AND 31
    ),
    CONSTRAINT chk_recorrencias_datas CHECK (
        data_fim IS NULL OR data_fim >= data_inicio
    )
);
\end{lstlisting}

\paragraph{Tabela: splitwise\_divisoes}
\begin{lstlisting}[caption=Estrutura da tabela splitwise_divisoes]
CREATE TABLE splitwise_divisoes (
    id UUID PRIMARY KEY DEFAULT gen_random_uuid(),
    transacao_id UUID NOT NULL REFERENCES transacoes(id) ON DELETE CASCADE,
    usuario_id UUID NOT NULL REFERENCES usuarios(id) ON DELETE CASCADE,
    valor_original DECIMAL(15,2) NOT NULL,
    percentual_divisao DECIMAL(5,2) DEFAULT 50.00,
    valor_devido DECIMAL(15,2) NOT NULL,
    pago BOOLEAN DEFAULT false,
    data_pagamento TIMESTAMP,
    observacoes TEXT,
    created_at TIMESTAMP DEFAULT CURRENT_TIMESTAMP,
    updated_at TIMESTAMP DEFAULT CURRENT_TIMESTAMP,
    
    CONSTRAINT chk_splitwise_percentual CHECK (
        percentual_divisao >= 0 AND percentual_divisao <= 100
    ),
    UNIQUE($transacao\_id$, $usuario\_id$)
);
\end{lstlisting}

\subsubsection{Tabelas de Auditoria}

\paragraph{Tabela: audit\_logs}
\begin{lstlisting}[caption=Estrutura da tabela audit_logs]
CREATE TABLE audit.audit_logs (
    id BIGSERIAL PRIMARY KEY,
    tabela VARCHAR(100) NOT NULL,
    operacao VARCHAR(10) NOT NULL,
    usuario_id UUID,
    dados_antigos JSONB,
    dados_novos JSONB,
    timestamp_operacao TIMESTAMP DEFAULT CURRENT_TIMESTAMP,
    ip_origem INET,
    user_agent TEXT,
    
    CONSTRAINT chk_audit_operacao CHECK (
        operacao IN ('INSERT', 'UPDATE', 'DELETE')
    )
) PARTITION BY RANGE ($timestamp\_operacao$);
\end{lstlisting}

\paragraph{Tabela: sessoes}
\begin{lstlisting}[caption=Estrutura da tabela sessoes]
CREATE TABLE sessoes (
    id UUID PRIMARY KEY DEFAULT gen_random_uuid(),
    usuario_id UUID NOT NULL REFERENCES usuarios(id) ON DELETE CASCADE,
    token_hash VARCHAR(255) NOT NULL,
    ip_origem INET NOT NULL,
    user_agent TEXT,
    data_criacao TIMESTAMP DEFAULT CURRENT_TIMESTAMP,
    data_expiracao TIMESTAMP NOT NULL,
    data_ultimo_acesso TIMESTAMP,
    ativa BOOLEAN DEFAULT true,
    
    CONSTRAINT chk_sessoes_expiracao CHECK (data_expiracao > data_criacao)
);
\end{lstlisting}

\subsection{Views Materializadas}

\paragraph{View: resumo\_financeiro\_mensal}
\begin{lstlisting}[caption=View materializada para resumos mensais]
CREATE MATERIALIZED VIEW resumo_financeiro_mensal AS
SELECT 
    usuario_id,
    grupo_id,
    DATE_TRUNC('month', data_transacao) AS mes_ano,
    SUM(CASE WHEN tipo = 'receita' THEN valor ELSE 0 END) AS total_receitas,
    SUM(CASE WHEN tipo = 'despesa' THEN valor ELSE 0 END) AS total_despesas,
    SUM(CASE WHEN tipo = 'receita' THEN valor ELSE -valor END) AS saldo_mensal,
    COUNT(*) AS total_transacoes
FROM transacoes 
WHERE status = 'confirmada'
GROUP BY usuario_id, grupo_id, DATE_TRUNC('month', data_transacao);

-- Indice para performance
CREATE UNIQUE INDEX idx_resumo_financeiro_pk 
    ON resumo_financeiro_mensal(usuario_id, COALESCE(grupo_id, '00000000-0000-0000-0000-000000000000'::uuid), mes_ano);
\end{lstlisting}

\subsection{Funcoes de Sistema}

\paragraph{Funcao: atualizar\_saldo\_conta}
\begin{lstlisting}[caption=Funcao para atualizar saldo das contas]
CREATE OR REPLACE FUNCTION atualizar_saldo_conta(conta_uuid UUID)
RETURNS DECIMAL(15,2) AS $$
DECLARE
    novo_saldo DECIMAL(15,2);
    saldo_inicial DECIMAL(15,2);
BEGIN
    -- Buscar saldo inicial
    SELECT c.saldo_inicial INTO saldo_inicial 
    FROM contas c WHERE c.id = conta_uuid;
    
    -- Calcular novo saldo
    SELECT saldo_inicial + COALESCE(SUM(
        CASE 
            WHEN t.tipo = 'receita' THEN t.valor
            WHEN t.tipo = 'despesa' THEN -t.valor
            WHEN t.tipo = 'transferencia' AND t.conta_id = conta_uuid THEN -t.valor
            WHEN t.tipo = 'transferencia' AND t.conta_destino_id = conta_uuid THEN t.valor
            ELSE 0
        END
    ), 0) INTO novo_saldo
    FROM transacoes t
    WHERE (t.conta_id = conta_uuid OR t.conta_destino_id = conta_uuid)
      AND t.status = 'confirmada';
    
    -- Atualizar saldo na conta
    UPDATE contas SET 
        saldo_atual = novo_saldo,
        updated_at = CURRENT_TIMESTAMP
    WHERE id = conta_uuid;
    
    RETURN novo_saldo;
END;
$$ LANGUAGE plpgsql;
\end{lstlisting}
